\cotefile{81}
\scanpage{35}
d'où
\[
\Sigma_{12}=X\deleted{\unsure{-}}\unsure{\setminus|-}(\Sigma_1^+\cup\Sigma_2^+)
\]
\unsure{donc} $\Sigma_{12}=\varnothing$ signifierait que
\[
X=\Sigma_1^+\amalg\Sigma_2^+
\]
et par suite
\[
\Pi_1=\left\{\Sigma_1^+,\underset{\inserted{\substack{\Vert\\\Sigma_{\unsure{1|2}}^-=X\setminus\Sigma_1^+}}}{\Sigma_2^+}\right\},
\quad \Pi_2=\{\Sigma_2^+,\Sigma_1^+\}
\]
alors $\deleted{\unsure{\Pi_2|\Pi_1}}\quad\Pi_1=\Pi_{\unsure{3|2}}$ contrairement \unsure{:|à} l'\unsure{hypothèse}.

On trouve donc

\noindent\makebox[0pt][r]{\rule[-2.7ex]{.4pt}{3.8ex}\hspace{.45em}}\deleted{\uline{Corollaire}} \inserted{Proposition} \unsure{Soient|Soit} $\Pi_1\unsure{;|,}\Pi_2$ deux décompositions binaires de $X$, mutuellement non \unsure{panachées}.
% Vertical emphasis strokes to the left of “binaires … Alors” and “tel”.
Alors il existe un \uline{unique} couple
\[
(\Sigma_1^{\unsure{+|\top}},\Sigma_2^+)\in\Pi_1\times\Pi_2
\]
\makebox[0pt][r]{\rule[-.35ex]{.4pt}{2ex}\hspace{.45em}}tel que l'on ait
\[
\Sigma_1^+\cap\Sigma_2^+=\varnothing
\]

\begingroup\makeatletter\tagsleft@true\makeatother
Soit alors
\[
C=\Pi_1\amalg\Pi_2=\unsure{\omega_1|w_1}\amalg\omega_2
\tag*{(17)}
\]
l'ensemble somme, qui a quatre éléments. On a une application canonique
\[
C\longrightarrow{\unsure{\mathfrak{P}|\mathcal{P}}}_{\unsure{\mathrm{finies}|\mathrm{ferm\acute{e}s}|\mathrm{foncier}}}(X)
\tag*{(1\deleted{\unsure{7}}8)}
\]
qui a comme image
\[
\{\Sigma_1^+,\Sigma_1^-,\Sigma_2^+,\Sigma_2^-\},
\tag*{(1\deleted{\unsure{8}}9)}
\]
\unsure{et|à} laquelle on l'identifie si $(1\deleted{\unsure{7}}8)$ est \unsure{injective}, i.e.

\endgroup
\clearpage
% archive sheet number: 34
\scanpage{36}
\begingroup\makeatletter\tagsleft@true\makeatother
si \deleted{\unsure{Card|le}} cardinal de l'ens. $(18)$ est $4$, i.e. si les quatre parties de $X$
\[
\Sigma_1^+,\quad\Sigma_1^-,\quad\Sigma_2^+,\quad\Sigma_2^-
\]
sont distinctes. A ce propos notons que l'on sait déjà par \deleted{\unsure{hypothèse} $(\Pi_1,\Pi_2$ \unsure{étant}} \inserted{définition} des déc. binaires
\[
\Sigma_1^+\ne\Sigma_1^-,\qquad\Sigma_2^+\ne\Sigma_2^-
\tag*{($*$)}
\]
et qu'on a
\[
\Sigma_1^+\ne\Sigma_2^+
\tag*{($**$)}
\]
puisque $\Sigma_1^+$ et $\unsure{\Sigma^+|\Sigma_2^+}$ sont des parties disjointes \uline{non} \uline{vides}. Notons donc qu'on lit sur le diagramme $(14)$
\[
\left\{\begin{aligned}
\text{a)}\;(\Sigma_1^+\ne\Sigma_2^-)&\Longleftrightarrow E_1\ne\Sigma_{12}
   &(\Longrightarrow E_2=\varnothing)\\
\text{b)}\;(\Sigma_2^+=\Sigma_1^-)&\unsure{\Longleftrightarrow|\Longrightarrow} E_2=\Sigma_{12}
   &(\Longrightarrow E_1=\varnothing)
\end{aligned}\right.
\tag*{($***$)}
\]
et que les deux relations \inserted{a), b)} \unsure{envisagées} dans $(***)$ ne peuvent être satisfaites simultanément, puisque $\deleted{\Sigma_{12}}$ l'on aurait $\Sigma_{12}=E_1=E_2=\varnothing$ puisque $E_1\cap E_2=\varnothing$, ce qui contredit $(16)$ (c).

Enfin
\[
(\Sigma_1^-=\Sigma_2^-)\Longleftrightarrow
(\Sigma_{\unsure{1\!2|1}}^-={\unsure{\Sigma|E}}_2^-=X)
\quad\text{i.e. }\Sigma_{12}=X
\tag*{($****$)}
\]
Ce cas est mutuellement exclusif avec chacun des cas $(***)$, \deleted{\unsure{dans ces derniers}} \inserted{\unsure{Car|car} si on avait en $=$ temps} $\deleted{\unsure{\Sigma_1^-}=\Sigma_2^-=X}\ \deleted{\unsure{E}}\Sigma_{12}=X$ et $E_1=\Sigma_{12}$, on aurait $E_1=X$ donc $\Sigma_1^+=\Sigma_1^-=X$, \deleted{\unsure{de|et}}\unsure{ce} qui contredit $(3)$.

\uline{Proposition} Pour que l'on ait $\operatorname{Card}\{\Sigma_1^+,\Sigma_1^-,\Sigma_2^+,\Sigma_{\unsure{2|1}}^-\}=4$, il faut et il suffit que l'on ne soit pas l'un des trois cas suivants, qui s'excluent mutuellement
\[
\left\{\begin{aligned}
\text{\unsure{a})}\quad &E_1=\Sigma_{12}&&(\Longrightarrow E_2=\varnothing)\\
\text{a')}\quad &E_2=\Sigma_{12}&&(\Longrightarrow E_1=\varnothing)\\
\text{b)}\quad &\Sigma_{12}=X\quad\text{i.e. }\Sigma_1^-=\Sigma_2^-=X
   &&(\text{\unsure{d'où|alors} }E_1=\Sigma_1^+,E_2=\Sigma_2^+)
\end{aligned}\right.
\tag*{(\deleted{\unsure{19|20}}20)}
\]
auquel cas le cardinal \unsure{du|des} $\unsure{\}\,4|\{\ \}|4}$ est égal à \unsure{trois}.

\endgroup
\clearpage
% archive sheet number: 35
\scanpage{37}
\begingroup\makeatletter\tagsleft@true\makeatother
\uline{\unsure{Cor|Cor.}} Dans le cas où $\operatorname{Card}\{\ \}=4$, \inserted{et $\Sigma_1^-\ne X,\Sigma_2^-\ne X$} les seules relations d'inclusions entre les ensembles $\Sigma_1^+,\Sigma_1^-,\Sigma_2^+,\Sigma_2^-$ sont les relations
% slanted
\margin{Sous forme : Pour que les relations \unsure{strictes} \unsure{(*)} soient les seules relations d'inclusion entre les ensembles $\Sigma_1^+,\Sigma_1^-,\Sigma_2^+,\Sigma_{\unsure{2|1}}^-$, il \unsure{f|f.} et s. que ceux-ci soient \inserted{\unsure{mutlt.}} distincts, et que l'on ait $\Sigma_1^-\ne X,\ \Sigma_2^-\ne X$.}
\[
\left\{\begin{aligned}
\Sigma_1^+&\subset\Sigma_2^-\\
\Sigma_2^+&\subset\Sigma_1^-
\end{aligned}\right. \quad .
\tag*{($*$)}
\]
En effet, on n'a pas les inclusions inverses, par hypothèse. D'autre part, on n'\uline{a pas}
\[
\Sigma_1^+\subset\Sigma_1^- \quad\text{ou}\quad \Sigma_1^-\unsure{\subset|\in}\Sigma_1^+,\quad\text{car on aurait}
\]
$\deleted{\unsure{\Sigma_1^+,\Sigma_1^+}}\ \Sigma_1^-=X$, de $=$ on n'a pas $\Sigma_2^+\subset{\unsure{\Sigma|E}}_{\unsure{2|1}}^-$ ni $\Sigma_2^-\subset\Sigma_2^+$, et on n'a pas $\Sigma_1^+\subset\Sigma_2^+$ ni l'inclusion inverse, puisque $\Sigma_1^+$ et $\unsure{\Sigma^+|\Sigma_2^+}$ sont disjoints et $\ne\varnothing$. Il reste à voir que l'on n'a pas une inclusion $\Sigma_1^-\subset\Sigma_2^-$ ou inclusion inverse, on lit sur le diagramme $(14)$ que \deleted{\unsure{cela}} \inserted{\unsure{la première|le premier} \unsure{p.ex.|p. ex.}} signifierait que $\deleted{\Sigma_{12}\unsure{\simeq|=}\Sigma_1^-}$ \deleted{\unsure{d'où}} $\Sigma_2^-=X$, ce que l'on a exclu.

% Vertical rule down the left of the definition.
\noindent\makebox[0pt][r]{\rule[-9ex]{.4pt}{10.5ex}\hspace{.55em}}\uline{\unsure{Déf|Déf.}} On dit que $\Pi_1,\Pi_2$ sont \uline{strictement non panachées} \inserted{si elles sont non panachées, et} si \deleted{\unsure{les}}les quatre ensembles $\Sigma_1^+,\Sigma_1^-,\Sigma_2^+,\Sigma_2^-$ sont distincts, \unsure{et.|i.e.} si \unsure{((|c} de plus $\Sigma_1^-\ne X,\Sigma_2^-\ne X$ -- i.e si les seules relations d'inclusion entre les ensembles $\Sigma_1^+,\Sigma_1^-,\Sigma_2^+,\Sigma_2^-$ sont $\Sigma_1^+\subset\Sigma_2^-,\ \Sigma_2^+\subset\Sigma_1^-$ -- ce qui signifie aussi que l'on a :
\[
\Sigma_1^-\ne X,\Sigma_2^-\ne X
\quad(\text{on dit \unsure{donc|alors} que }\Pi_1,\Pi_2\text{ sont \uline{propres}}),
\ E_1\ne\Sigma_{12},E_2\ne\Sigma_{12}.
\tag*{(21)}
\]

\endgroup
\clearpage
% archive sheet number: 36
\scanpage{38}
\begingroup\makeatletter\tagsleft@true\makeatother
Mais si \unsure{cette|ces} condition n'est pas satisfaite pour $\deleted{\unsure{\Sigma|E}}\Pi_1,\Pi_2$ non panachées, on reste en \unsure{plan}, et on regarde l'ens. à quatre éléments
\[
C=\Pi_1\amalg\Pi_2=\unsure{\omega_1|w_1}\amalg\omega_2
=\{{\unsure{\varepsilon|\epsilon}}_1^+,\varepsilon_1^-,\varepsilon_2^+,\varepsilon_2^-\}
\tag*{(17)}
\]
et la relation d'ordre dessus, où les seules relations $<$ sont données par
\[
\varepsilon_1^+<\varepsilon_2^-,\qquad\varepsilon_2^+<\varepsilon_1^-
\tag*{(22)}
\]
\diagram{\begin{tikzcd}[ampersand replacement=\&,row sep=2.3em,column sep=2.6em]
|[label=left:{\(\varepsilon_1^-\)}]| \bullet \arrow[from=2-2,to=1-1]
 \& |[label=right:{\(\varepsilon_2^-\)}]| \bullet\\
|[label=left:{\(\varepsilon_1^+\)},label=below:{\(\deleted{\unsure{\varepsilon_1^+|\varepsilon_2^+}}\)}]| \bullet \arrow[ur,crossing over]
 \& |[label=right:{\(\varepsilon_2^+\)}]| \bullet
\end{tikzcd}}

1.3. Soit \deleted{$(\Pi_\alpha)_{\alpha\in A}$ une famille}
\[
(\Pi_\alpha)_{\alpha\in A},\qquad \Pi_\alpha=\{\Sigma_\alpha^\varepsilon\}_{\varepsilon\in\omega_\alpha}
\tag*{(23)}
\]
% slanted
\margin{c'est un élément de $\unsure{\mathfrak{P}|\mathcal{P}}(\mathfrak{P}_2(\deleted{\unsure{X|\times}}\unsure{\mathfrak{P}|\mathcal{P}|\square})(X))$}
\unsure{une|un} \deleted{famille} \inserted{\unsure{ens.} (\unsure{écrit|écrite} \unsure{indiciellement|individuellement})} \unsure{de|des} décompositions binaires de $X$, on suppose \deleted{qu'elles}
\[
\forall\alpha,\beta\in A,\quad\alpha\ne\beta,\qquad
\Pi_\alpha,\Pi_\beta\ \text{sont mut\textsuperscript{t}. non panachées}.
\tag*{(24)}
\]
% slanted
\margin{On dit, si $\operatorname{Card} A=\unsure{n|\nu}$, que l'on a une \uline{décomposition} \uline{$n$-aire} non panachée de $X$}

\inserted{\uline{\unsure{Déf|Déf.}}} On dit \deleted{alors} que $(\Pi_\alpha)_{\alpha\in A}$ est non \unsure{panachée|panachées}. (Terminologie analogue pour \uline{strictement} non panachée).

Considérons
\[
C=\coprod_{\alpha\in A}\omega_\alpha
\tag*{(25)}
\]
\inserted{avec l'involution ss pt fixe $\sigma$ qui, sur chaque $\omega_\alpha$, induit une \deleted{\illegible{1}} transposition}
et sur $C$ la relation d'ordre, dans laquelle on a

\endgroup
\clearpage
% archive sheet number: 37
\scanpage{39}
\begingroup\makeatletter\tagsleft@true\makeatother
\[
\left(\underset{\substack{\unsure{\rotatebox{-90}{$\in$}|\rotatebox{-90}{$\ni$}}\\\unsure{\omega_\alpha|w_\alpha}}}{x}
<
\underset{\substack{\rotatebox{-90}{$\in$}\\\omega_\beta}}{y}\right)
\Longleftrightarrow
\Biggl(\alpha\ne\beta,\quad
\underbrace{x\underset{\{\Pi_\alpha,\Pi_\beta\}}{<}y}_{
  \inserted{\substack{\text{i.e. }\deleted{\unsure{x}}\\\Sigma_\alpha^x\cap\Sigma_\beta^{\sigma y}=\varnothing}}}
\Biggr)
\tag*{(26)}
\]
% slanted
\margin{il faut vérifier la \unsure{transitivité} de cette relation -- c'est immédiat}
Alors $\sigma$ est un antiautomorphisme (involutif sans pt fixe) de $C$. L'application
\[
\begin{aligned}
C&\xrightarrow{\unsure{\varphi|\phi}}{\unsure{\mathfrak{P}|\mathcal{P}}}_{\unsure{\mathrm{finies}|\mathrm{ferm\acute{e}s}|\mathrm{foncier}}}(X)\\
(\alpha,x)&\longmapsto\Sigma_\alpha^x
\end{aligned}
\tag*{(27)}
\]
est une application croissante. Pour qu'\unsure{elles|elle} soit injective, et définisse un isomorphisme d'ens ordonnés de $C$ avec son image, il f et s. que $(\Pi_\alpha)_{\alpha\in A}$ soit strictement non \unsure{panaché|panachée}.

\unsure{Indépendamment} de $X$ :

\uline{Définition} On appelle \uline{\unsure{maquette|moquette} \deleted{\unsure{de}} \unsure{ordonnée}} (de \unsure{décompositions} \unsure{$n$-aires|$n$ aires}) un ens \deleted{\unsure{ordonné}} \inserted{fini} $C$, muni d'une involution \inserted{$\sigma$} sans pt fixe, et d'une relation d'ordre, ayant la \unsure{propriété|propriétés} \unsure{suivante|suivantes} :

\deleted{\unsure{a)} Si $x,y\in C$ et $x<y$, alors $x$ et $y$}
\deleted{Si $x,y\in C$, $\unsure{\sigma x\ne y|\sigma x=y}$, $x$ et $y$ non \unsure{congrus} sous $\sigma$, alors on n'a pas $x<y$ non \unsure{congrus} \unsure{mod} $C$,}
\deleted{\unsure{b)}} \deleted{\unsure{et}} \deleted{\unsure{b)}}
% slanted
\margin{NB implique que $x,y$ \unsure{compar.} $\Longrightarrow$ $x,y$ non \unsure{congrus} mod $\sigma$}
\[
\left\{\begin{aligned}
\text{a)}\;&\text{Si }x,y\in C,\quad\deleted{\unsure{x\ne y}},\quad
\text{alors }\deleted{\sigma x,y\text{ et }x,\sigma y}\ 
\inserted{x,y\text{ comparables}}\\
&\text{sss }\ \sigma x,y\ \text{non comparables}\\
\text{b)}\;&\text{Si }x,y\in C,\ \deleted{\text{comparables, et }\sigma y<\sigma x}\
\text{ alors }\ x<y\Longrightarrow\sigma y<\sigma x
\end{aligned}\right.
\tag*{(\unsure{28})}
\]
\uline{NB} Cela \unsure{indique|implique} que $\sigma$ est un antiautomorphisme de l'ens ordonné $C$

\endgroup
\clearpage
% archive sheet number: 38
\scanpage{40}
\noindent\makebox[0pt][r]{\rule[-1.4ex]{.4pt}{3.7ex}\hspace{.5em}}\uline{Proposition} Soi\deleted{\unsure{ent|t}}t $(C,\sigma,<)$ une \unsure{maquette|moquette} ordonnée. Alors pour tt $x\in C$, les \deleted{\unsure{ens.|sous}} ens.
\[
\begin{aligned}
 C_{\geq x}&=\{y\in C\mid y\geq x\}\\
 C_{\leq x}&=\{y\in C\mid y\leq x\}
\end{aligned}
\]
\makebox[0pt][r]{\rule[-.4ex]{.4pt}{2.4ex}\hspace{.5em}}sont \uline{tot\textsuperscript{\unsure{t|l}}} ordonnés.

Il suffit de \inserted{le} prouver \inserted{pour} le premier \deleted{\unsure{notons|noter}} (par dualité), \deleted{Soit \unsure{donc}} \unsure{cela|cel} signifie aussi :
\[
\begin{aligned}
 &\text{Si}\quad x<y,\quad x<z,\quad \text{avec}\quad y\ne z\\
 &\hspace{1em}\text{alors}\quad (y,z)\ \text{comparables.}
\end{aligned}
\]

\noindent\begin{minipage}[c]{.38\linewidth}
\diagram{\begin{tikzcd}[ampersand replacement=\&,row sep={7.4mm,between origins},column sep={12.4mm,between origins},
execute at end picture={
\draw[line width=.6pt] ([xshift=-14mm,yshift=-1.1mm]Jfortyx.center) -- ([xshift=-10mm,yshift=.2mm]Jfortyx.center);
\draw[line width=.45pt] ([xshift=-3.4mm,yshift=2.3mm]Jfortyx.center) -- ([xshift=-1.7mm,yshift=2.8mm]Jfortyx.center);
\draw[->,line width=.6pt] ([xshift=4mm,yshift=4.2mm]Jfortyx.center) -- ([xshift=8.6mm,yshift=6.3mm]Jfortyx.center);
\draw[line width=.75pt] ([xshift=14mm,yshift=-.8mm]Jfortyx.center) -- ([xshift=6.3mm,yshift=1.9mm]Jfortyx.center);
\draw[->,line width=.85pt] ([xshift=1.2mm,yshift=3.1mm]Jfortyx.center) -- ([xshift=-8mm,yshift=6.2mm]Jfortyx.center);
}]
 |[inner sep=0pt,label=above left:{\(z\)}]| \phantom{\cdot} \& \cdot \& |[inner sep=0pt,label=right:{\(y\)}]| \cdot\\
 \& |[alias=Jfortyx,inner sep=0pt,label=below:{\(x\)}]| \cdot\arrow[ul,shorten <=2.9mm,shorten >=.3mm]\arrow[ur,shorten <=2mm,shorten >=.5mm] \&
\end{tikzcd}}
\end{minipage}\hfill\begin{minipage}[c]{.57\linewidth}
[On note que $y$ et $z$ sont non congrus mod $\sigma$, à cause de 28 a)\unsure{,|.}]
\end{minipage}

Si $\deleted{\unsure{x}}y$ et $\deleted{\unsure{y}}z$ étaient incomparables, on aurait que $\sigma y$ et $z$ sont comparables, On \unsure{peu|pas|ne} peut avoir $z\unsure{<|\leq}\sigma y$ car on\unsure{.} aurait $x<\sigma y$, ce qui contredit $x<y$ par (28) a). On \deleted{ne peut avoir} \inserted{aurait donc} $\sigma y<z$ d'où $\sigma z<\unsure{\backslash|x}$

\clearpage
% archive sheet number: 39
\scanpage{41}
\begin{flushright}
\unsure{relations|relatons} d'\unsure{equiv.|èquiv.|équiv.|equi.}\hspace*{.8em}\\
non \unsure{panachées|panachés}
\end{flushright}

% archive sheet number: 40
